\chapter{Round surfaces and circular Willmore tori}\label{sec:examples-willmore-tori}
These calculations use the averaged energy $\mathcal W=\frac14\int H^2$ in Euclidean space and $\mathcal W_1=\int(1+H^2/4)$ in the unit three-sphere, where $H$ is the trace. The circular torus and its relation to the Clifford torus appear in \cite{MNWill}, p.~684; the calculations here supply their metrics, curvature and exact energies. The global positive-genus theorem remains the separate Theorem \ref{thm:willmore-bound}.

\begin{proposition}[Geometry of a circular ring torus]\label{prop:circular-ring-torus-geometry}
For $R>r>0$, the map
\[
 T(\theta,\varphi)=\bigl((R+r\cos\theta)\cos\varphi,
 (R+r\cos\theta)\sin\varphi,r\sin\theta\bigr),
 \qquad (\theta,\varphi)\in(\mathbb R/2\pi\mathbb Z)^2,
\]
is an embedded torus with metric $r^2d\theta^2+(R+r\cos\theta)^2d\varphi^2$ and area $4\pi^2Rr$. With the outward normal
$\nu=(\cos\theta\cos\varphi,\cos\theta\sin\varphi,\sin\theta)$,
its principal curvatures are $1/r$ and $\cos\theta/(R+r\cos\theta)$.
\end{proposition}
\begin{proof}
Differentiate the parametrization. The two coordinate derivatives are orthogonal and have respective lengths $r$ and $R+r\cos\theta>0$, so it is an immersion with the stated metric. The positive distance to the rotation axis determines $\varphi$, while that distance minus $R$ together with height determines $(r\cos\theta,r\sin\theta)$ and hence $\theta$. This proves injectivity and thus embeddedness by compactness. The area density is $r(R+r\cos\theta)$, whose double integral is $4\pi^2Rr$. Differentiating the stated normal gives $\nu_\theta=T_\theta/r$ and $\nu_\varphi=\cos\theta\,T_\varphi/(R+r\cos\theta)$, proving the principal curvatures with $A=D\nu$.
\end{proof}

\begin{theorem}[The energy minimum within the circular ring family]\label{thm:circular-ring-willmore-energy}
The preceding torus has
\[
 \mathcal W(T)=\frac{\pi^2R^2}{r\sqrt{R^2-r^2}}\geq2\pi^2,
\]
with equality exactly when $R/r=\sqrt2$. Its energy tends to infinity as $R/r\downarrow1$ or $R/r\to\infty$. This is the circular-family calculation motivating the conjecture recalled in \cite{LY82}, p.~270, and \cite{MNWill}, p.~684; it does not prove the inequality for arbitrary immersed tori.
\end{theorem}
\begin{proof}
Proposition \ref{prop:circular-ring-torus-geometry} gives
\[
 \mathcal W(T)=\frac{\pi}{2r}\int_0^{2\pi}
 \frac{(R+2r\cos\theta)^2}{R+r\cos\theta}\,d\theta.
\]
Polynomial division makes the integrand $4r\cos\theta+R^2/(R+r\cos\theta)$. The cosine integrates to zero. The substitution $t=\tan(\theta/2)$ on $(-\pi,\pi)$ gives
\[
 \int_0^{2\pi}\frac{d\theta}{R+r\cos\theta}
 =2\int_{-\infty}^{\infty}\frac{dt}{(R+r)+(R-r)t^2}
 =\frac{2\pi}{\sqrt{R^2-r^2}}.
\]
This proves the formula. Put $a=R/r>1$. The derivative of $a^2/\sqrt{a^2-1}$ is $a(a^2-2)/(a^2-1)^{3/2}$, negative below $\sqrt2$ and positive above it. Its value there is two, and the endpoint limits are infinite. Only the scale-free ratio matters, as conformal dilation invariance predicts.
\end{proof}

\begin{example}[Total Gauss curvature does not determine the torus energy]\label{ex:ring-torus-total-curvature}
Every circular ring torus above satisfies
\[
 \int_TK\,d\mu=0,\qquad \int_T|K|\,d\mu=8\pi,
 \qquad \int_{\{K>0\}}K\,d\mu=4\pi,
\]
although its Willmore energy varies from $2\pi^2$ to arbitrarily large values.
\end{example}
\begin{proof}
The principal curvatures give $K=\cos\theta/[r(R+r\cos\theta)]$. Multiplying by area leaves $\cos\theta\,d\theta d\varphi$. The integrals of $\cos\theta$, $|\cos\theta|$ and its positive part over a period are zero, four and two. Multiply each by $2\pi$. The energy variation follows from Theorem \ref{thm:circular-ring-willmore-energy}. Thus even fixing these total intrinsic-curvature quantities does not fix the bending energy.
\end{proof}

\begin{proposition}[Product tori in the unit three-sphere]\label{prop:unit-sphere-product-tori-energy}
For $a,b>0$ with $a^2+b^2=1$, the product torus
\[
 F(\theta,\varphi)=(a\cos\theta,a\sin\theta,b\cos\varphi,b\sin\varphi)
\]
has area $4\pi^2ab$, intrinsic Gauss curvature zero, and spherical principal curvatures $b/a$ and $-a/b$. Its spherical Willmore energy is $\pi^2/(ab)\geq2\pi^2$, with equality exactly at $a=b=1/\sqrt2$. Within this family, only that Clifford torus is a Willmore critical surface.
\end{proposition}
\begin{proof}
The product metric is $a^2d\theta^2+b^2d\varphi^2$, flat with the claimed area. The field
$\nu=(b\cos\theta,b\sin\theta,-a\cos\varphi,-a\sin\varphi)$
is a unit vector tangent to $S^3$ and normal to the torus. Its derivatives in the two product directions give $b/a$ and $-a/b$. Thus $H=(b^2-a^2)/(ab)$ is constant, and
\[
 1+\frac{H^2}{4}=\frac{4a^2b^2+(b^2-a^2)^2}{4a^2b^2}
 =\frac1{4a^2b^2}.
\]
Multiplication by area gives the energy. The inequality $ab\leq(a^2+b^2)/2=1/2$ proves its minimum. The trace-free curvature is nonzero everywhere, because the two principal curvatures have opposite signs. Hence the equation of Theorem \ref{thm:willmore-first-variation-spaceform} reduces to $H=0$, or $a=b$. The sphere here has radius one; it is different from the radius-two sphere containing the Clifford self-shrinker of Example \ref{ex:clifford-torus-shrinker}.
\end{proof}

\begin{example}[Stereographic projection gives the circular minimizer]\label{ex:clifford-stereographic-ring}
The stereographic projection from $(0,0,0,1)$ of the unit Clifford torus is the ring torus with major radius $\sqrt2$ and minor radius one. Both descriptions therefore have Willmore energy $2\pi^2$.
\end{example}
\begin{proof}
Write a Clifford point as $(\cos\varphi,\sin\varphi,\cos\theta,\sin\theta)/\sqrt2$. Its projected cylindrical coordinates are
\[
 \rho=\frac1{\sqrt2-\sin\theta},\qquad
 z=\frac{\cos\theta}{\sqrt2-\sin\theta}.
\]
The denominator is positive. Direct algebra gives $\rho^2+z^2+1=2\sqrt2\rho$, or $(\rho-\sqrt2)^2+z^2=1$. Thus the image lies in the stated ring torus. Stereographic projection is a diffeomorphism off its pole, and the Clifford torus avoids the pole, so its restriction is an embedding. Its image in the ring torus is open by the inverse function theorem and closed by compactness; connectedness gives the entire ring torus. Proposition \ref{prop:unit-sphere-product-tori-energy}, Theorem \ref{thm:circular-ring-willmore-energy} and the stereographic energy identity all give the same value.
\end{proof}

\begin{example}[All round latitudes have the same Willmore energy]\label{ex:latitude-willmore-energy}
For $0<\alpha<\pi$, the round latitude $F(\omega)=(\sin\alpha\,\omega,\cos\alpha)$ in $S^3$, $\omega\in S^2$, has area $4\pi\sin^2\alpha$, trace mean curvature $2\cot\alpha$ for a suitable normal, and $\mathcal W_1=4\pi$. Every latitude is Willmore critical, although only the equator is minimal.
\end{example}
\begin{proof}
The induced metric is $\sin^2\alpha$ times the unit-sphere metric. The unit normal $\nu=(\cos\alpha\,\omega,-\sin\alpha)$ differentiates to $A=\cot\alpha I$, proving the curvature and minimality assertion. The energy is $4\pi\sin^2\alpha(1+\cot^2\alpha)=4\pi$. Since $A^0=0$ and $H$ is constant, the Willmore equation vanishes. The ambient-curvature term one is necessary: the integral of $H^2/4$ alone is $4\pi\cos^2\alpha$ and is not constant in this spherical family.
\end{proof}

\begin{example}[An affine plane becomes a punctured round sphere]\label{ex:inversion-affine-plane-sphere}
In a three-dimensional affine subspace identified with $\mathbb R^3$, invert at zero the plane $\{x\cdot n=d\}$, where $|n|=1$ and $d\ne0$. Its image is the sphere of center $n/(2d)$ and radius $1/(2|d|)$ with the origin removed. The plane has zero Willmore energy, while the punctured sphere has energy $4\pi$.
\end{example}
\begin{proof}
Writing $z=x/|x|^2$, the plane equation becomes $z\cdot n=d|z|^2$, equivalently
$|z-n/(2d)|^2=1/(4d^2)$. Every nonzero point of this sphere inverts to a point of the plane, proving the image assertion. A plane is minimal; the round sphere's energy is $4\pi$, and removing one point does not change its integral. The original plane is noncompact, so invariance for a closed surface with a pole off its image does not apply in that direction. Inverting the completed sphere at its origin instead fits Proposition \ref{prop:willmore-inversion-energy-defect}: it loses exactly $4\pi$ and becomes the plane.
\end{proof}

\begin{proposition}[The full Willmore equation on a ring torus]\label{prop:ring-torus-willmore-equation}
For the circular ring torus of Proposition \ref{prop:circular-ring-torus-geometry}, write $s=R+r\cos\theta$. Then
\[
 \Delta H+H|A^0|^2=\frac{R(R^2-2r^2)}{2r^3s^3}.
\]
Consequently the ratio $R/r=\sqrt2$ gives a critical point under all smooth variations, not merely under variations within the ring family.
\end{proposition}
\begin{proof}
The trace is $H=1/r+\cos\theta/s$, and its derivative is $H_\theta=-R\sin\theta/s^2$. For a function of $\theta$ alone, the product metric gives
\[
 \Delta H=\frac1{r^2}H_{\theta\theta}+\frac{s_\theta}{r^2s}H_\theta
 =-\frac{R\cos\theta}{r^2s^2}-\frac{R\sin^2\theta}{rs^3}.
\]
Also $|A^0|^2=(1/r-\cos\theta/s)^2/2=R^2/(2r^2s^2)$. Bringing the sum to the common denominator $2r^3s^3$ leaves
\[
 R\bigl[-2r\cos\theta\,s-2r^2\sin^2\theta+R(s+r\cos\theta)\bigr]
 =R(R^2-2r^2).
\]
The denominator is everywhere positive. The first-variation theorem therefore gives the assertion, including the converse within this circular family. A minimum found in a restricted family would not by itself justify the full criticality conclusion; the displayed equation does.
\end{proof}

\begin{proposition}[Second variation at a round Euclidean sphere]\label{prop:round-sphere-willmore-hessian}
For the round sphere $S_R^2\subset\mathbb R^3$, a smooth variation with initial normal speed $f$ has
\[
 \left.\frac{d^2}{dt^2}\right|_0\mathcal W(F_t)
 =\frac12\int_{S_R^2}(\Delta f)^2d\mu-\frac1{R^2}\int_{S_R^2}|\nabla f|^2d\mu
 =\int_{S_R^2}\left|\nabla^2f-\frac{\Delta f}{2}g\right|^2d\mu.
\]
In particular every such second variation is nonnegative, independently of the second-order choice of the variation.
\end{proposition}
\begin{proof}
On the sphere, $H=2/R$ is constant and $A^0=0$. Differentiate the integrand of Theorem \ref{thm:willmore-first-variation-spaceform}. Since the original Euler--Lagrange expression is zero, derivatives of speed and area do not contribute. The derivative of $H|A^0|^2$ is zero. The derivative of the Laplacian acting on the original constant $H$ is also zero, because every Laplacian annihilates a constant. Thus only $\Delta\dot H$ remains. The Euclidean linearization gives $\dot H=-\Delta f-2f/R^2$. This yields
\[
 \ddot{\mathcal W}=\frac12\int f\Delta(\Delta+2/R^2)f,
\]
which is the first formula after integration by parts.

For completeness, commuting one covariant derivative of $\nabla f$ gives $\operatorname{div}(\nabla^2f)=\nabla\Delta f+\operatorname{Ric}(\nabla f,\cdot)$. On this two-sphere $\operatorname{Ric}=R^{-2}g$. Integrating the divergence identity against $\nabla f$ gives
\[
 \int|\nabla^2f|^2=\int(\Delta f)^2-\frac1{R^2}\int|\nabla f|^2.
\]
Subtracting half the squared trace produces exactly the preceding Hessian formula. Nonnegativity is now the integral of a squared norm; no spectral decomposition is needed.
\end{proof}

\begin{proposition}[Only dilations and translations lie in the spherical kernel]\label{prop:round-sphere-willmore-kernel}
The second variation in Proposition \ref{prop:round-sphere-willmore-hessian} vanishes exactly for functions $f=c+a\cdot X$, with $c\in\mathbb R$ and $a\in\mathbb R^3$. The constant mode is the normal speed of a change of radius; the linear modes are normal translation speeds up to a fixed factor. The normal kernel is therefore four-dimensional.
\end{proposition}
\begin{proof}
Vanishing of the square gives $\nabla^2f=(\Delta f/2)g$. Take its divergence and use the commutation identity from the preceding proof. It follows that $\nabla(\Delta f+2f/R^2)=0$. Averaging this constant expression over the sphere shows that $u=f-\overline f$ satisfies $\Delta u=-2u/R^2$ and hence $\nabla^2u=-ug/R^2$.

Regard $\nabla u$ as an ambient tangent vector. The Euclidean derivative of
$Z=\nabla u+(u/R^2)X$ in a tangent direction $V$ is zero: its tangential terms are $-uV/R^2+uV/R^2$, and its normal terms are $-(V u)X/R^2+(V u)X/R^2$, because the sphere's vector-valued second fundamental form is $B(V,W)=-g(V,W)X/R^2$. Connectedness makes $Z=a$ one constant ambient vector. Dotting with $X$ gives $u=a\cdot X$. Conversely every such linear function has Hessian $-ug/R^2$, and constants have zero Hessian, so they lie in the kernel. A translation by a vector $b$ has normal speed $b\cdot X/R$, which spans the same three-dimensional space. These are infinitesimal kernel statements and do not assert convergence of a geometric flow.
\end{proof}
